% !TEX root = ../main.tex
\begin{ouverture}[piles]{8}{Piles et oxydoréduction}{Partie 1 · Constitution et transformations de la matière}
{Oxydants et réducteurs\enspace·\enspace Transfert d'électrons\enspace·\enspace Fonctionnement d'une pile\enspace·\enspace Capacité électrique}
\objectif{Citer des oxydants et des réducteurs usuels ; justifier le caractère réducteur des métaux du bloc s.}
\objectif{Illustrer un transfert spontané d'électrons par contact entre réactifs ou par un circuit extérieur.}
\objectif{Justifier la séparation des réactifs dans deux demi-piles et le rôle du pont salin.}
\objectif{Modéliser et schématiser le fonctionnement d'une pile à partir de résultats expérimentaux.}
\objectif{Déterminer la capacité électrique d'une pile à partir de sa constitution initiale.}
\prerequis{Oxydant, réducteur, couple, demi-équation électronique (classe de première).}
\prerequis{Quotient de réaction, constante d'équilibre, critère d'évolution (chapitre~7).}
\prerequis{Intensité du courant, tension, charge électrique $Q=I\,\Delta t$.}
\end{ouverture}
\essentiel
\essSection{1}{Oxydants et réducteurs}
\begin{definition}[Couple oxydant-réducteur]
Un \textbf{oxydant} est une espèce capable de capter un ou plusieurs électrons ; un \textbf{réducteur} est une espèce
capable d'en céder. Un oxydant Ox et un réducteur Red forment un couple Ox/Red lorsqu'ils sont liés par une
demi-équation électronique :
\[\mathrm{Ox}+n\,\mathrm{e}^-\rightleftharpoons\mathrm{Red}.\]
Une \textbf{oxydation} est une perte d'électrons, une \textbf{réduction} un gain d'électrons. Une réaction
d'oxydoréduction est un transfert d'électrons du réducteur d'un couple vers l'oxydant d'un autre couple.
\end{definition}
\begin{savoirfaire}[Écrire l'équation d'une réaction d'oxydoréduction]
Écrire les deux demi-équations dans le sens de la transformation : équilibrer les éléments autres que O et H, puis
l'oxygène avec \ce{H2O}, l'hydrogène avec \ce{H+}, enfin les charges avec des électrons. Multiplier chaque demi-équation
par un coefficient pour que les électrons cédés et captés soient en même nombre, puis additionner : les électrons
n'apparaissent pas dans l'équation de la réaction.
\end{savoirfaire}
\begin{exemple}
L'acide ascorbique (vitamine~C) réduit le diiode :
\begin{center}\renewcommand{\arraystretch}{1.25}
\begin{tabular}{@{}r@{\quad}l@{}}
réduction : & \ce{I2 + 2 e- -> 2 I-}\\
oxydation : & \ce{C6H8O6 -> C6H6O6 + 2 H+ + 2 e-}\\\midrule
bilan : & \ce{I2 + C6H8O6 -> 2 I- + C6H6O6 + 2 H+}
\end{tabular}\end{center}
Cette réaction rapide et totale permet de doser la vitamine~C d'un jus de fruit.
\end{exemple}
\begin{propriete}[Caractère réducteur des métaux du bloc s]
Les métaux du bloc s (lithium, sodium, potassium, magnésium, calcium\,\dots) n'ont qu'un ou deux électrons de valence,
qu'ils cèdent facilement pour acquérir la configuration électronique du gaz noble qui les précède : \ce{Na -> Na+ + e-},
\ce{Mg -> Mg^2+ + 2 e-}. Ce sont des réducteurs puissants : le sodium réagit vivement avec l'eau en libérant du
dihydrogène,
\[\ce{2 Na(s) + 2 H2O(l) -> 2 Na^+(aq) + 2 HO^-(aq) + H2(g)}.\]
On les conserve à l'abri de l'eau et de l'air.
\end{propriete}
\begin{propriete}[Oxydants et réducteurs usuels]
\par\smallskip
{\small\renewcommand{\arraystretch}{1.3}
\begin{tabularx}{\linewidth}{@{}lllY@{}}\toprule
\otTete{Espèce} & \otTete{Couple} & \otTete{Demi-équation électronique} & \otTete{Usage}\\\midrule
\multicolumn{4}{@{}l}{\textit{Oxydants}}\\
dioxygène & \ce{O2}/\ce{H2O} & \ce{O2 + 4 H+ + 4 e- <=> 2 H2O} & respiration, combustions\\
dichlore & \ce{Cl2}/\ce{Cl-} & \ce{Cl2 + 2 e- <=> 2 Cl-} & traitement de l'eau\\
eau de Javel (\ce{ClO-}) & \ce{ClO-}/\ce{Cl-} & \ce{ClO- + 2 H+ + 2 e- <=> Cl- + H2O} & désinfection\\\midrule
\multicolumn{4}{@{}l}{\textit{Réducteurs}}\\
dihydrogène & \ce{H+}/\ce{H2} & \ce{2 H+ + 2 e- <=> H2} & pile à combustible\\
acide ascorbique & \ce{C6H6O6}/\ce{C6H8O6} & \ce{C6H6O6 + 2 H+ + 2 e- <=> C6H8O6} & antioxydant (E300)\\
métal M & \ce{M^{$n$+}}/\ce{M} & \ce{M^{$n$+} + $n$ e- <=> M} & piles, corrosion\\\bottomrule
\end{tabularx}}
\end{propriete}
\essSection{2}{Transfert spontané d'électrons}
\begin{propriete}[Transfert direct ou par un circuit extérieur]
Une lame de zinc plongée dans une solution de sulfate de cuivre(II) se recouvre de cuivre et la solution se décolore :
les ions \ce{Cu^2+} captent directement les électrons cédés par le zinc,
\[\ce{Cu^2+(aq) + Zn(s) -> Cu(s) + Zn^2+(aq)}.\]
Ce transfert spontané respecte le critère d'évolution : $Q_{\mathrm{r},\mathrm{i}}=[\ce{Zn^2+}]_\mathrm{i}/[\ce{Cu^2+}]_\mathrm{i}$ est très
inférieur à $K=\num{1.9e37}$. L'énergie libérée est perdue sous forme de chaleur. Si l'on sépare les deux couples dans
deux compartiments reliés par un fil conducteur, les électrons transitent par le fil : on obtient un courant électrique
utilisable.
\end{propriete}
\schema{08-contact}
\essSection{3}{Constitution et fonctionnement d'une pile}
\begin{definition}[Pile électrochimique]
Une pile est constituée de deux \textbf{demi-piles}, formées chacune d'une électrode (conducteur électronique, souvent
un métal) plongeant dans une solution qui contient l'autre espèce du couple. Les deux solutions sont reliées par un
\textbf{pont salin} (solution ionique gélifiée, par exemple de nitrate de potassium) ou par une membrane. La
\textbf{tension à vide} est la tension entre les bornes de la pile lorsqu'elle ne débite pas ; elle se mesure avec un
voltmètre.
\end{definition}
\begin{propriete}[Fonctionnement d'une pile]
Lorsque la pile débite, la réaction d'oxydoréduction évolue spontanément dans le sens prévu par le critère d'évolution
($Q_{\mathrm{r},\mathrm{i}}<K$ : sens direct) :
\begin{itemize}[leftmargin=1.4em,itemsep=0.1em]
\item à l'\textbf{anode}, pôle négatif, le réducteur est oxydé et libère des électrons dans le circuit ;
\item à la \textbf{cathode}, pôle positif, l'oxydant capte ces électrons et est réduit ;
\item dans le circuit extérieur, les électrons circulent du pôle $-$ vers le pôle $+$ ; le courant circule en sens inverse ;
\item dans les solutions et le pont salin, le courant est dû au déplacement des ions : les cations migrent vers la
cathode, les anions vers l'anode. Le pont salin ferme le circuit et maintient l'électroneutralité des solutions.
\end{itemize}
\end{propriete}
\schema{08-pile}
\begin{savoirfaire}[Déterminer la polarité et le fonctionnement d'une pile]
Brancher un voltmètre aux bornes de la pile : si la tension affichée est positive, la borne V est reliée au pôle $+$ et la
borne COM au pôle $-$. Écrire la réduction à l'électrode positive, l'oxydation à l'électrode négative, puis l'équation de
fonctionnement en combinant les deux demi-équations.
\end{savoirfaire}
\begin{attention}
Les électrons ne circulent ni dans les solutions ni dans le pont salin : seuls les ions s'y déplacent. Séparer les
réactifs est indispensable : mélangés, ils échangeraient directement leurs électrons, sans courant utilisable.
\end{attention}
\begin{experience}[Réaliser une pile Daniell]
Une lame de zinc décapée plonge dans 50~mL de solution de sulfate de zinc à \qty{0,10}{mol.L^{-1}}, une lame de cuivre
dans 50~mL de solution de sulfate de cuivre(II) à \qty{0,10}{mol.L^{-1}} ; une bande de papier filtre imbibée de solution
de nitrate de potassium relie les deux béchers. Un voltmètre dont la borne COM est reliée au zinc affiche \qty{1,10}{V} :
le zinc est le pôle négatif. Sans le pont salin, la tension mesurée est nulle. Si la pile débite plusieurs heures dans
un conducteur ohmique, la lame de cuivre s'épaissit, la solution bleue pâlit et la lame de zinc s'amincit.
\end{experience}
\essSection{4}{Usure et capacité d'une pile}
\begin{propriete}[Usure d'une pile]
Au cours du fonctionnement, le quotient de réaction augmente. La pile cesse de débiter lorsque $Q_\mathrm{r}$ atteint $K$
ou, comme c'est le cas pour la plupart des piles usuelles dont la constante est très grande, lorsqu'un réactif est
épuisé : la pile est usée.
\end{propriete}
\begin{definition}[Capacité électrique d'une pile]
La capacité électrique $Q$ d'une pile est la charge électrique maximale qu'elle peut faire circuler :
\[Q=n(\mathrm{e}^-)_{\max}\times F,\qquad F=N_\mathrm{A}\,e=\qty{96500}{C.mol^{-1}},\]
où $n(\mathrm{e}^-)_{\max}$ est la quantité d'électrons échangés jusqu'à l'épuisement du réactif limitant et $F$ la
constante de Faraday. Elle s'exprime en coulombs ou en ampères-heures ($\qty{1}{A.h}=\qty{3600}{C}$). Une pile de
capacité $Q$ qui débite un courant d'intensité constante $I$ fonctionne pendant $\Delta t=Q/I$.
\end{definition}
\begin{savoirfaire}[Déterminer la capacité d'une pile]
Écrire l'équation de fonctionnement et repérer le nombre $n$ d'électrons échangés pour un avancement $x$ :
$n(\mathrm{e}^-)=n\,x$. Déterminer le réactif limitant et l'avancement maximal $x_{\max}$, puis calculer $Q=n\,x_{\max}\,F$.
\end{savoirfaire}
\begin{exemple}
La pile Daniell précédente contient \qty{5,0e-3}{mol} d'ions \ce{Cu^2+} (50~mL à \qty{0,10}{mol.L^{-1}}) et une lame de
zinc de 2,0~g, soit \qty{3,1e-2}{mol} : les ions \ce{Cu^2+} sont limitants. Deux électrons sont échangés par ion \ce{Cu^2+}
réduit, d'où $Q=2\times\num{5.0e-3}\times\num{96500}=\qty{965}{C}$, soit \qty{0,27}{A.h}. En débitant \qty{20}{mA}, la pile
fonctionne pendant $\Delta t=965/0{,}020=\qty{4,8e4}{s}$, soit 13~h environ.
\end{exemple}
\begin{codecote}[0.58]
\begin{pythoncol}
F = 96500                    # C/mol
def capacite(reactifs, n):
    # reactifs : (quantité, coefficient)
    x_max = min(q / a for q, a in reactifs)
    return n * x_max * F     # en C
Q = capacite([(5.0e-3, 1), (3.1e-2, 1)], 2)
print(round(Q), "C")
print(round(Q / 3600, 2), "A.h")
print(round(Q / 0.020 / 3600, 1), "h")
\end{pythoncol}
\begin{consolecol}
965 C
0.27 A.h
13.4 h
\end{consolecol}
\cote
\fontsize{9.5}{12.5}\selectfont
\capacite{Python}\par\smallskip
La fonction \code{capacite} reçoit les quantités initiales des réactifs consommés et leurs nombres stœchiométriques.
L'avancement maximal est le plus petit des quotients quantité/coefficient : il désigne le réactif limitant. La
capacité vaut alors $n\,x_{\max}\,F$. On retrouve les valeurs de l'exemple.
\end{codecote}
\schema{08-usure}
\begin{rediger}
Pour décrire une pile, préciser pour chaque électrode le pôle, la nature de la transformation et la demi-équation écrite
dans le sens où elle se produit : « à l'électrode de cuivre, pôle positif, les ions cuivre(II) sont réduits :
\ce{Cu^2+ + 2 e- -> Cu} ». Pour une capacité, justifier le réactif limitant et donner le résultat en coulombs et en
ampères-heures.
\end{rediger}
\exercices

\rubrique{Découverte}
\begin{blocLivre}{ex-8-1}
\begin{exercice}[1]{Écrire des demi-équations électroniques}
On considère les couples oxydant-réducteur suivants.
\begin{center}\renewcommand{\arraystretch}{1.3}
\begin{tabular}{@{}l@{\hspace{12mm}}l@{\hspace{12mm}}l@{}}
\otLabel{A.}~\ce{Ag+}/\ce{Ag} & \otLabel{B.}~\ce{Fe^3+}/\ce{Fe^2+} & \otLabel{C.}~\ce{I2}/\ce{I-}\\
\otLabel{D.}~\ce{O2}/\ce{H2O} & \otLabel{E.}~\ce{H+}/\ce{H2} & \otLabel{F.}~\ce{MnO4-}/\ce{Mn^2+}
\end{tabular}\end{center}
\begin{questions}
\item Écrire la demi-équation électronique de chaque couple.
\item Dans la réaction \ce{Cu(s) + 2 Ag^+(aq) -> Cu^2+(aq) + 2 Ag(s)}, quelle espèce est oxydée ? quelle espèce est
réduite ? Quel est l'oxydant ? le réducteur ?
\item Écrire l'équation de la réaction entre les ions fer(II) et le dioxygène en milieu acide.
\item Écrire l'équation de la réaction entre les ions permanganate et les ions fer(II) en milieu acide.
\end{questions}
\end{exercice}
\end{blocLivre}
\begin{blocLivre}{ex-8-2}
\begin{exercice}[1]{L'arbre de Diane}
On plonge un fil de cuivre dans 50,0~mL d'une solution de nitrate d'argent à \qty{0,100}{mol.L^{-1}}. Au bout de quelques
heures, de fins cristaux d'argent ramifiés recouvrent le fil et la solution bleuit. La réaction
\ce{Cu(s) + 2 Ag^+(aq) <=> Cu^2+(aq) + 2 Ag(s)} a pour constante $K\approx\num{4e15}$ à 25~°C.
\schema{08-ex02}
\begin{donnees}
$M(\ce{Ag})=\qty{107,9}{g.mol^{-1}}$ ; $M(\ce{Cu})=\qty{63,5}{g.mol^{-1}}$.
\end{donnees}
\begin{questions}
\item Quelles observations montrent la formation d'argent métallique ? d'ions cuivre(II) ?
\item Écrire les deux demi-équations électroniques. Identifier l'oxydant et le réducteur.
\item Calculer le quotient de réaction initial et justifier que la transformation est spontanée.
\item Le cuivre est en excès. Calculer la masse maximale d'argent qui se dépose et la masse de cuivre consommée.
\item Peut-on récupérer de l'énergie électrique avec ce dispositif ? Justifier.
\end{questions}
\end{exercice}
\end{blocLivre}
\begin{blocLivre}{ex-8-3}
\begin{exercice}[1]{Déterminer la polarité d'une pile}
Une lame de cuivre plonge dans une solution de nitrate de cuivre(II), une lame d'argent dans une solution de nitrate
d'argent ; les deux solutions, à \qty{0,10}{mol.L^{-1}}, sont reliées par un pont salin au nitrate de potassium. Un
voltmètre est branché comme sur le schéma.
\schema{08-ex03}
\begin{questions}
\item Déterminer la polarité de la pile. Justifier.
\item Écrire les réactions aux électrodes lorsque la pile débite. Préciser l'anode et la cathode.
\item Écrire l'équation de fonctionnement. Montrer qu'elle est en accord avec le critère d'évolution, sachant que
$K\approx\num{4e15}$.
\item On remplace le voltmètre par un conducteur ohmique. Indiquer sur un schéma le sens de déplacement des électrons,
le sens du courant et le sens de déplacement des ions du pont salin.
\end{questions}
\end{exercice}
\end{blocLivre}
\begin{blocLivre}{ex-8-4}
\begin{exercice}[1]{Le rôle du pont salin}
Une pile Daniell (zinc dans une solution de sulfate de zinc, cuivre dans une solution de sulfate de cuivre(II), pont
salin au nitrate de potassium) alimente un conducteur ohmique en série avec un ampèremètre. On réalise trois essais.
\begin{questions}[label=\otLabel{\alph*.},leftmargin=1.9em,itemsep=0pt,parsep=0.15em]
\item Avec le pont salin, l'ampèremètre indique \qty{12}{mA}.
\item Sans le pont salin, l'ampèremètre indique 0.
\item On verse la solution de sulfate de cuivre(II) dans le bécher qui contient la lame de zinc : un dépôt rougeâtre
apparaît sur le zinc, mais l'ampèremètre indique 0.
\end{questions}
\begin{questions}
\item Quels porteurs de charge assurent le passage du courant dans les fils ? dans les solutions et le pont salin ?
\item Interpréter l'essai b.
\item Dans quel sens les ions \ce{K+} et \ce{NO3-} du pont salin se déplacent-ils lorsque la pile débite ? Justifier à
l'aide de l'électroneutralité des solutions.
\item Interpréter l'essai c. Pourquoi faut-il séparer les réactifs dans une pile ?
\end{questions}
\end{exercice}
\end{blocLivre}
\begin{blocLivre}{ex-8-5}
\begin{exercice}[1]{Oxydants et réducteurs du quotidien}
\begin{questions}[label=\otLabel{\alph*.},leftmargin=1.9em,itemsep=0pt,parsep=0.15em]
\item L'eau de Javel, solution d'ions hypochlorite \ce{ClO-}, désinfecte les surfaces.
\item Le jus de citron, riche en acide ascorbique \ce{C6H8O6}, empêche une pomme coupée de brunir à l'air.
\item Le dioxygène de l'air fait rouiller le fer.
\item Le dihydrogène alimente les piles à combustible.
\end{questions}
\begin{questions}
\item Pour chaque situation, identifier l'espèce citée et préciser s'il s'agit d'un oxydant ou d'un réducteur.
\item Écrire les demi-équations des couples \ce{ClO-}/\ce{Cl-}, \ce{C6H6O6}/\ce{C6H8O6} et \ce{O2}/\ce{H2O}.
\item Expliquer pourquoi la vitamine~C protège la pomme.
\item L'eau de Javel contient aussi des ions chlorure. En milieu acide, il se produit la réaction
\ce{ClO-(aq) + Cl-(aq) + 2 H+(aq) -> Cl2(g) + H2O(l)}. Identifier l'oxydant et le réducteur. Pourquoi ne faut-il jamais
mélanger l'eau de Javel avec un détartrant acide ?
\end{questions}
\end{exercice}
\end{blocLivre}
\rubrique{Entraînement}
\begin{blocLivre}{ex-8-6}
\begin{exercice}[2]{La pile Daniell}
Une lame de zinc de masse 5,0~g plonge dans 100~mL de solution de sulfate de zinc à \qty{0,10}{mol.L^{-1}} ; une lame de
cuivre plonge dans 100~mL de solution de sulfate de cuivre(II) à \qty{0,10}{mol.L^{-1}}. Les solutions sont reliées par un
pont salin. Un voltmètre dont la borne COM est reliée au zinc affiche \qty{1,10}{V}.
\begin{donnees}
\ce{Cu^2+(aq) + Zn(s) <=> Cu(s) + Zn^2+(aq)} : $K=\num{1.9e37}$ ; $M(\ce{Zn})=\qty{65,4}{g.mol^{-1}}$ ;
$M(\ce{Cu})=\qty{63,5}{g.mol^{-1}}$ ; $F=\qty{96500}{C.mol^{-1}}$.
\end{donnees}
\begin{questions}
\item Quelle est la polarité de la pile ? Écrire les demi-équations des réactions aux électrodes.
\item Calculer le quotient de réaction initial. Le sens d'évolution est-il en accord avec la polarité observée ?
\item Déterminer le réactif limitant, puis la capacité de la pile en coulombs et en milliampères-heures.
\item La pile débite un courant d'intensité constante \qty{50}{mA}. Pendant combien de temps peut-elle fonctionner ?
\item De quelle masse la lame de cuivre s'est-elle alourdie lorsque la pile est usée ?
\end{questions}
\end{exercice}
\end{blocLivre}
\begin{blocLivre}{ex-8-7}
\begin{exercice}[2]{La pile d'une montre}
Une pile bouton à l'oxyde d'argent alimente une montre. Elle contient 34~mg de zinc en poudre et 0,15~g d'oxyde
d'argent \ce{Ag2O}, séparés par une membrane imbibée d'une solution concentrée d'hydroxyde de potassium. Sa tension vaut
1,55~V. Les réactions aux électrodes sont :
\[\ce{Zn + 2 HO- -> ZnO + H2O + 2 e-}\qquad\text{et}\qquad\ce{Ag2O + H2O + 2 e- -> 2 Ag + 2 HO-}.\]
\begin{donnees}
$M(\ce{Zn})=\qty{65,4}{g.mol^{-1}}$ ; $M(\ce{Ag2O})=\qty{231,7}{g.mol^{-1}}$ ; $F=\qty{96500}{C.mol^{-1}}$ ;
\mbox{1 an $=\qty{3,156e7}{s}$}.
\end{donnees}
\begin{questions}
\item Quelle électrode constitue le pôle positif ? Justifier.
\item Écrire l'équation de fonctionnement de la pile. Que remarque-t-on pour les ions hydroxyde ?
\item Déterminer le réactif limitant, puis la capacité de la pile en milliampères-heures.
\item La montre consomme en moyenne un courant de \qty{1,5}{\micro A}. Estimer la durée de fonctionnement de la pile.
\item Pourquoi ces piles usagées doivent-elles être collectées ?
\end{questions}
\end{exercice}
\end{blocLivre}
\begin{blocLivre}{ex-8-8}
\begin{exercice}[2]{Suivre l'usure d'une pile}
Une pile Daniell est formée d'une lame de zinc, en excès, dans 100~mL de solution de sulfate de zinc à
\qty{0,100}{mol.L^{-1}} et d'une lame de cuivre dans 100~mL de solution de sulfate de cuivre(II) à \qty{0,100}{mol.L^{-1}}.
Elle débite un courant d'intensité constante $I$. Une sonde de spectrophotométrie permet de suivre la concentration des
ions \ce{Cu^2+}.
\schema{08-ex08}
\begin{donnees}
$K=\num{1.9e37}$ pour \ce{Cu^2+(aq) + Zn(s) <=> Cu(s) + Zn^2+(aq)} ; $F=\qty{96500}{C.mol^{-1}}$.
\end{donnees}
\begin{questions}
\item Écrire l'équation de fonctionnement de la pile.
\item Montrer que la concentration des ions \ce{Cu^2+} diminue linéairement au cours du temps. Déterminer le
coefficient directeur de la droite, en \unit{mol.L^{-1}.h^{-1}}.
\item Exprimer la quantité d'électrons qui a circulé entre les dates 0 et $t$ en fonction de la variation de
$[\ce{Cu^2+}]$. En déduire l'intensité $I$ du courant.
\item Au bout de combien de temps la pile sera-t-elle usée ? Calculer sa capacité.
\item À cet instant, le quotient de réaction a-t-il atteint la constante d'équilibre ? Pourquoi la pile s'arrête-t-elle ?
\end{questions}
\end{exercice}
\end{blocLivre}
\begin{blocLivre}{ex-8-9}
\begin{exercice}[2]{La pile zinc-air des aides auditives}
Les aides auditives sont alimentées par des piles bouton « zinc-air » : la taille 675 a une tension de 1,4~V et une
capacité de \qty{620}{mA.h}. À l'achat, un opercule autocollant bouche les trous d'aération ; une fois l'opercule retiré,
le dioxygène de l'air pénètre dans la pile. Les réactions aux électrodes sont :
\[\ce{Zn + 2 HO- -> ZnO + H2O + 2 e-}\qquad\text{et}\qquad\ce{O2 + 2 H2O + 4 e- -> 4 HO-}.\]
\begin{donnees}
$M(\ce{Zn})=\qty{65,4}{g.mol^{-1}}$ ; $F=\qty{96500}{C.mol^{-1}}$.
\end{donnees}
\begin{questions}
\item Identifier l'oxydant et le réducteur. Quelle électrode constitue le pôle négatif ?
\item Écrire l'équation de fonctionnement de la pile.
\item Convertir la capacité en coulombs, puis calculer la masse minimale de zinc contenue dans la pile.
\item Une aide auditive consomme en moyenne \qty{1,6}{mA}, 16~heures par jour. Combien de jours la pile dure-t-elle ?
\item Pourquoi ces piles sont-elles légères pour leur capacité ? Pourquoi ne retire-t-on l'opercule qu'au dernier moment ?
\end{questions}
\end{exercice}
\end{blocLivre}
\begin{blocLivre}{ex-8-10}
\begin{exercice}[2]{Le sodium, un métal très réducteur}
Le sodium, métal du bloc s ($Z=11$), est conservé sous de l'huile. Derrière un écran de protection, le professeur dépose
un morceau de sodium de masse 0,23~g sur l'eau d'un cristallisoir additionnée de phénolphtaléine : le métal fond, se
déplace en dégageant un gaz, et une coloration rose apparaît.
\begin{donnees}
$M(\ce{Na})=\qty{23,0}{g.mol^{-1}}$ ; volume molaire des gaz : $V_\mathrm{m}=\qty{24}{L.mol^{-1}}$ ; la phénolphtaléine est
rose si $\pH>8{,}2$ ; couple \ce{H2O}/\ce{H2} : \ce{2 H2O + 2 e- <=> H2 + 2 HO-}.
\end{donnees}
\begin{questions}
\item Écrire la configuration électronique de l'atome de sodium. Combien possède-t-il d'électrons de valence ?
\item Justifier que le sodium est un réducteur puissant. Quel ion forme-t-il ?
\item Le gaz dégagé est du dihydrogène. Écrire l'équation de la réaction et interpréter la coloration rose.
\item Calculer le volume de dihydrogène dégagé.
\item Pourquoi conserve-t-on le sodium sous de l'huile ? Écrire la demi-équation du couple \ce{Mg^2+}/\ce{Mg} : en quoi le
magnésium ($Z=12$), lui aussi du bloc s, ressemble-t-il au sodium ?
\end{questions}
\end{exercice}
\end{blocLivre}
\begin{blocLivre}{ex-8-11}
\begin{exercice}[2]{Simuler la décharge d'une pile}
Une pile Daniell contient \qty{5,0e-3}{mol} d'ions \ce{Cu^2+} et \qty{5,0e-3}{mol} d'ions \ce{Zn^2+} ; le zinc est en excès.
Le programme suivant simule son fonctionnement lorsqu'elle débite un courant d'intensité constante.
\begin{python}
F, I, dt = 96500, 0.050, 60      # C/mol, A, s
n_cu, n_zn = 5.0e-3, 5.0e-3      # mol
t = 0
while n_cu > 0:
    dx = I * dt / (2 * F)        # avancement pendant dt
    n_cu = n_cu - dx
    n_zn = n_zn + dx
    t = t + dt
print(t, round(t / 3600, 2), round(n_zn, 4))
\end{python}
L'exécution affiche \code{19320 5.37 0.01}.
\begin{questions}
\item \capacite{Python}Écrire l'équation de fonctionnement de la pile. Que représentent les variables \code{n\_cu} et
\code{n\_zn} ?
\item Justifier l'expression de \code{dx}.
\item Pourquoi la boucle s'arrête-t-elle ? Interpréter les trois nombres affichés.
\item Retrouver la durée de fonctionnement à partir de la capacité de la pile. Pourquoi le programme donne-t-il une
valeur légèrement supérieure ?
\item Quelle hypothèse du modèle n'est plus réaliste en fin de décharge ?
\end{questions}
\end{exercice}
\end{blocLivre}
\rubrique{Objectif bac}
\begin{blocLivre}{ex-8-12}
\begin{exercice}[3]{La pile d'un stimulateur cardiaque}
\begin{docu}{1}{Une pile implantée}
Un stimulateur cardiaque délivre au cœur des impulsions électriques. Depuis les années 1970, il est alimenté par une
pile lithium-iode : l'anode est en lithium, la cathode en diiode mélangé à un polymère conducteur. Entre les deux se
forme une fine couche d'iodure de lithium solide, qui joue le rôle d'électrolyte. La pile ne contient aucune solution
aqueuse ; sa tension vaut 2,8~V.
\end{docu}
\begin{docu}{2}{Durée de vie}
Une pile de stimulateur a une capacité d'environ \qty{1,0}{A.h} ; le stimulateur consomme en moyenne un courant de
\qty{10}{\micro A}. Lorsque la pile est usée, le boîtier est remplacé au cours d'une intervention chirurgicale.
\end{docu}
\schema{08-ex12}
\begin{donnees}
Lithium : $Z=3$, $M=\qty{6,94}{g.mol^{-1}}$ ; $M(\ce{I})=\qty{126,9}{g.mol^{-1}}$ ; $M(\ce{Zn})=\qty{65,4}{g.mol^{-1}}$ ;
$F=\qty{96500}{C.mol^{-1}}$ ; \mbox{1 an $=\qty{3,156e7}{s}$}.
\end{donnees}
\begin{questions}
\item Le lithium est un métal du bloc s. Justifier son caractère réducteur. Écrire les demi-équations des couples
\ce{Li+}/\ce{Li} et \ce{I2}/\ce{I-}.
\item Préciser l'anode, la cathode et la polarité de la pile, puis écrire l'équation de fonctionnement.
\item Pourquoi la pile ne contient-elle aucune solution aqueuse ?
\item \initiative Estimer la durée de fonctionnement du stimulateur, puis les masses minimales de lithium et de diiode
contenues dans la pile.
\item Quelle masse de zinc faudrait-il pour obtenir la même capacité avec une anode en zinc ? Conclure sur l'intérêt du
lithium.
\end{questions}
\end{exercice}
\end{blocLivre}
\begin{blocLivre}{ex-8-13}
\begin{exercice}[3]{Une voiture à hydrogène}
\begin{docu}{1}{La pile à combustible}
La Toyota Mirai de deuxième génération emporte 5,6~kg de dihydrogène, comprimé à 700~bar dans trois réservoirs ; son
autonomie annoncée est d'environ 650~km. Elle est propulsée par un moteur électrique alimenté par une pile à combustible
qui ne rejette que de l'eau.
\source{D'après Toyota, présentation de la Mirai, 2020.}
\end{docu}
\begin{docu}{2}{Une cellule de la pile}
Chaque cellule comporte deux électrodes séparées par une membrane qui ne laisse passer que les ions \ce{H+}. Les
réactions aux électrodes sont :
\[\ce{H2 -> 2 H+ + 2 e-}\qquad\text{et}\qquad\ce{O2 + 4 H+ + 4 e- -> 2 H2O}.\]
En fonctionnement, la tension aux bornes d'une cellule vaut environ 0,70~V.
\end{docu}
\schema{08-ex13}
\begin{donnees}
$M(\ce{H2})=\qty{2,0}{g.mol^{-1}}$ ; $M(\ce{H2O})=\qty{18,0}{g.mol^{-1}}$ ; $F=\qty{96500}{C.mol^{-1}}$ ; énergie électrique
fournie par une charge $Q$ sous la tension $U$ : $E=Q\,U$ ; \mbox{$\qty{1}{kW.h}=\qty{3,6e6}{J}$}.
\end{donnees}
\begin{questions}
\item Laquelle des électrodes A et B est l'anode ? la cathode ? Préciser le sens de déplacement des électrons dans le
circuit extérieur et celui des ions \ce{H+} dans la membrane.
\item Écrire l'équation de fonctionnement. À quel élément d'une pile de laboratoire la membrane correspond-elle ?
\item Calculer la charge électrique que peut fournir le dihydrogène d'un plein.
\item \initiative Estimer l'énergie électrique disponible, puis la consommation de la voiture en \unit{kW.h} pour 100~km.
Comparer à une voiture électrique à batterie, qui consomme environ \qty{15}{kW.h} pour 100~km.
\item Quelle masse d'eau la voiture rejette-t-elle en consommant un plein de dihydrogène ?
\end{questions}
\end{exercice}
\end{blocLivre}
\begin{blocLivre}{ex-8-14}
\begin{exercice}[3]{Une pile de concentration}
\begin{docu}{1}{Deux demi-piles presque identiques}
Deux lames de cuivre de 2,0~g plongent chacune dans 100~mL de solution de sulfate de cuivre(II) : à
\qty{1,0e-2}{mol.L^{-1}} dans le compartiment 1, à \qty{0,10}{mol.L^{-1}} dans le compartiment 2. Les solutions sont
reliées par un pont salin. Un voltmètre dont la borne COM est reliée à la lame du compartiment 1 affiche
$U=+\qty{0,030}{V}$.
\end{docu}
\schema{08-ex14}
\begin{donnees}
$M(\ce{Cu})=\qty{63,5}{g.mol^{-1}}$ ; $F=\qty{96500}{C.mol^{-1}}$.
\end{donnees}
\begin{questions}
\item Déterminer la polarité de la pile, puis écrire les réactions aux électrodes.
\item Écrire l'équation de fonctionnement et l'expression de son quotient de réaction en fonction des concentrations
$[\ce{Cu^2+}]_1$ et $[\ce{Cu^2+}]_2$. Justifier que $K=1$.
\item Vérifier que le sens d'évolution prévu par le critère d'évolution est en accord avec la polarité observée.
\item \initiative Déterminer la composition des deux solutions lorsque la pile est usée, puis la capacité de la pile.
La lame du compartiment 1 est-elle assez massive ?
\item En quoi l'usure de cette pile diffère-t-elle de celle d'une pile Daniell ?
\end{questions}
\end{exercice}
\end{blocLivre}
\begin{blocLivre}{ex-8-15}
\begin{exercice}[3]{Classer des métaux à l'aide de piles}
\begin{docu}{1}{Trois piles}
On réalise des demi-piles en plongeant une lame de métal dans une solution de ses ions à \qty{0,10}{mol.L^{-1}}. On
associe deux demi-piles par un pont salin et on mesure la tension à vide.
\par\smallskip
{\small\renewcommand{\arraystretch}{1.2}
\begin{tabular}{@{}lcc@{}}\toprule
Pile & Borne COM reliée à & Tension affichée\\\midrule
zinc -- cuivre & la lame de zinc & $+\qty{1,10}{V}$\\
cuivre -- argent & la lame d'argent & $-\qty{0,43}{V}$\\
fer -- cuivre & la lame de cuivre & $-\qty{0,78}{V}$\\\bottomrule
\end{tabular}}
\end{docu}
\begin{docu}{2}{Un clou cuivré}
Un clou en fer plongé dans une solution de sulfate de cuivre(II) se recouvre d'un dépôt de cuivre.
\end{docu}
\begin{questions}
\item Déterminer la polarité de chacune des trois piles.
\item Pour la pile zinc -- cuivre, écrire les réactions aux électrodes et l'équation de fonctionnement.
\item On admet que les tensions à vide s'ajoutent comme des écarts d'altitude : sur une échelle verticale, le pôle
positif d'une pile est placé au-dessus du pôle négatif, à une distance égale à la tension à vide. Placer le zinc, le fer
et l'argent sur cette échelle, le cuivre étant à l'origine.
\item \initiative On admet aussi que le métal le plus réducteur constitue le pôle négatif d'une pile. Classer les quatre
métaux par pouvoir réducteur croissant, puis prévoir la tension à vide et la polarité des piles zinc -- argent et
zinc -- fer.
\item Le document~2 est-il en accord avec ce classement ? Écrire l'équation de la réaction.
\end{questions}
\end{exercice}
\end{blocLivre}
